\documentclass[10pt,a4paper]{amsart}
\usepackage{amsmath,amssymb,amsthm,graphicx}
% Retained minimal preamble for the Stacks Project chapter more-algebra.tex (More on Algebra).
% The chapter file pinned-original-source.tex carries only \input{preamble} on its first line, so no
% preamble text lives inside the chapter itself.  The Stacks Project uses one single shared preamble for
% every chapter; the snapshot stages it identically as sources/stacks-chow/preamble.tex and
% sources/stacks-etale-cohomology/preamble.tex.  The block below is copied mechanically, line for line,
% from sources/stacks-chow/preamble.tex lines 14-18 (the Xy-pic setup) and 154-217 (the native theorem
% declarations together with the mathematical macros).  Nothing was added, removed or redefined; the
% wrapper class and the amsmath/amssymb/amsthm/graphicx base packages are supplied by the assembler.
% Chapter: more-algebra.tex -- More on Algebra.
% For commutative diagrams we use Xy-pic
\usepackage[all]{xy}

% We use 2cell for 2-commutative diagrams.
\xyoption{2cell}

% Theorem environments.
%
\theoremstyle{plain}
\newtheorem{theorem}[subsection]{Theorem}
\newtheorem{proposition}[subsection]{Proposition}
\newtheorem{lemma}[subsection]{Lemma}

\theoremstyle{definition}
\newtheorem{definition}[subsection]{Definition}
\newtheorem{example}[subsection]{Example}
\newtheorem{exercise}[subsection]{Exercise}
\newtheorem{situation}[subsection]{Situation}

\theoremstyle{remark}
\newtheorem{remark}[subsection]{Remark}
\newtheorem{remarks}[subsection]{Remarks}

\numberwithin{equation}{subsection}

% Macros
%
\def\lim{\mathop{\mathrm{lim}}\nolimits}
\def\colim{\mathop{\mathrm{colim}}\nolimits}
\def\Spec{\mathop{\mathrm{Spec}}}
\def\Hom{\mathop{\mathrm{Hom}}\nolimits}
\def\Ext{\mathop{\mathrm{Ext}}\nolimits}
\def\SheafHom{\mathop{\mathcal{H}\!\mathit{om}}\nolimits}
\def\SheafExt{\mathop{\mathcal{E}\!\mathit{xt}}\nolimits}
\def\Sch{\mathit{Sch}}
\def\Mor{\mathop{\mathrm{Mor}}\nolimits}
\def\Ob{\mathop{\mathrm{Ob}}\nolimits}
\def\Sh{\mathop{\mathit{Sh}}\nolimits}
\def\NL{\mathop{N\!L}\nolimits}
\def\CH{\mathop{\mathrm{CH}}\nolimits}
\def\proetale{{pro\text{-}\acute{e}tale}}
\def\etale{{\acute{e}tale}}
\def\QCoh{\mathit{QCoh}}
\def\Ker{\mathop{\mathrm{Ker}}}
\def\Im{\mathop{\mathrm{Im}}}
\def\Coker{\mathop{\mathrm{Coker}}}
\def\Coim{\mathop{\mathrm{Coim}}}

% Boxtimes
%
\DeclareMathSymbol{\boxtimes}{\mathbin}{AMSa}{"02}

%
% Macros for moduli stacks/spaces
%
\def\QCohstack{\mathcal{QC}\!\mathit{oh}}
\def\Cohstack{\mathcal{C}\!\mathit{oh}}
\def\Spacesstack{\mathcal{S}\!\mathit{paces}}
\def\Quotfunctor{\mathrm{Quot}}
\def\Hilbfunctor{\mathrm{Hilb}}
\def\Curvesstack{\mathcal{C}\!\mathit{urves}}
\def\Polarizedstack{\mathcal{P}\!\mathit{olarized}}
\def\Complexesstack{\mathcal{C}\!\mathit{omplexes}}
% \Pic is the operator that assigns to X its picard group, usage \Pic(X)
% \Picardstack_{X/B} denotes the Picard stack of X over B
% \Picardfunctor_{X/B} denotes the Picard functor of X over B
\def\Pic{\mathop{\mathrm{Pic}}\nolimits}
\def\Picardstack{\mathcal{P}\!\mathit{ic}}
\def\Picardfunctor{\mathrm{Pic}}
\def\Deformationcategory{\mathcal{D}\!\mathit{ef}}

% --- appended 2026-10-12 (Chief): the project own definitions of the reference/slogan/history
% display environments, copied mechanically from the shared Stacks preamble lines 8-12
% (\usepackage{verbatim} plus the three \newenvironment lines, which the chapter preamble keeps
% verbatim from the project preamble). Nothing added, removed or redefined: these environments
% carry no mathematical content and the project itself typesets them as comments. Without these
% three lines a counted statement whose own text uses them fails with Environment slogan undefined.
\usepackage{verbatim}
\newenvironment{reference}{\comment}{\endcomment}
\newenvironment{slogan}{\comment}{\endcomment}
\newenvironment{history}{\comment}{\endcomment}
\usepackage{cleveref}
\crefname{originallemma}{original lemma}{original lemmas}
\Crefname{originallemma}{Original lemma}{Original lemmas}
\makeatletter
\newlabel{more-algebra-lemma-limit-flat}{{0912}{source}{}{sourcepointer-1}{}}
\newlabel{more-algebra-lemma-limit-flat@cref}{{[originallemma][0][]0912}{[1][0][]source}{}{}{}}
\makeatother
\crefname{originallemma}{original lemma}{original lemmas}
\Crefname{originallemma}{Original lemma}{Original lemmas}
\makeatletter
\newlabel{lemma-tor-strictly-pro-zero}{{0911}{source}{}{sourcepointer-2}{}}
\newlabel{lemma-tor-strictly-pro-zero@cref}{{[originallemma][0][]0911}{[1][0][]source}{}{}{}}
\makeatother
\crefname{originallemma}{original lemma}{original lemmas}
\Crefname{originallemma}{Original lemma}{Original lemmas}
\makeatletter
\newlabel{homology-lemma-apply-Mittag-Leffler-again}{{070E}{source}{}{sourcepointer-3}{}}
\newlabel{homology-lemma-apply-Mittag-Leffler-again@cref}{{[originallemma][0][]070E}{[1][0][]source}{}{}{}}
\makeatother
\crefname{originallemma}{original lemma}{original lemmas}
\Crefname{originallemma}{Original lemma}{Original lemmas}
\makeatletter
\newlabel{algebra-lemma-flat}{{00HD}{source}{}{sourcepointer-4}{}}
\newlabel{algebra-lemma-flat@cref}{{[originallemma][0][]00HD}{[1][0][]source}{}{}{}}
\makeatother
\title{Flatness of an inverse limit of flat modules over the quotients by powers of an ideal}
\author{The Stacks Project Authors}
\date{stacks-project:more-algebra.tex:local-snapshot-20261012; GNU FDL 1.2 or later}
\begin{document}
\maketitle
\setcounter{subsection}{1}
\noindent\textit{Editorial note.} All delivered mathematical text is the author's own, copied from the pinned Stacks Project chapter file more-algebra.tex; no statement or proof step was rewritten, reordered or removed.\par\medskip
\noindent\textit{Editorial note.} The chapter file pinned-original-source.tex carries only \textbackslash{}input\{preamble\} on its first line, so no preamble text lives inside the chapter. The retained file sources/stacks-more-algebra/retained-preamble.tex is a mechanical, line-for-line copy of the project's single shared preamble as staged in the snapshot (Xy-pic setup plus theorem declarations and mathematical macros). Nothing was added, removed or redefined, and the excerpt serves only as preamble, never as a span coordinate.\par\medskip
\noindent\textit{Editorial note.} Both counted spans start at the \textbackslash{}begin line of their environment, so the statement and the author's proof are delivered whole; surrounding prose of the section is deliberately outside the counted spans and is carried as standard prerequisites instead.\par\medskip
\noindent\textit{Editorial note.} The citation fields (source\_id, source\_version, authors, source\_url, license.name, source\_title) are identical to those of every other card this lane has produced from this chapter, so the batch selector treats them as one source.\par\medskip
\noindent\textit{Allowed standard prerequisites (editorial).}\begin{itemize}
\item Tor against a quotient by a power of an ideal is an essentially constant pro-zero system. Lemma \textbackslash{}ref\{lemma-tor-strictly-pro-zero\}, cited at proof lines 6498-6499. Makes the obstruction to commuting tensor products with the limit vanish.
\item Mittag-Leffler style lemma identifying the limit of a cokernel with the cokernel of the limits. Homology, Lemma \textbackslash{}ref\{homology-lemma-apply-Mittag-Leffler-again\}, cited at proof lines 6500 and 6519. Produces the isomorphism between the tensor product with the limit and the limit of the tensor products, and the injectivity of the limit of the maps.
\item Flatness is detected by injectivity of tensoring with injective maps of modules. Algebra, Lemma \textbackslash{}ref\{algebra-lemma-flat\}, cited at proof lines 6508-6509. Turns the second part of the proof into a criterion for flatness of the limit module.
\end{itemize}\medskip
\ifdefined\hypertarget\hypertarget{sourcepointer-1}{}\fi
\noindent\textit{Source reference.} Original lemma 0912: focused-seeds/AGP/sources/stacks-snapshot-20261011/tags.txt, source lines 11430--11430. The lemma's own canonical tag, quoted so that readers can resolve the label inside the delivered text; it points at the counted statement itself and carries no proof content.\par\medskip
\ifdefined\hypertarget\hypertarget{sourcepointer-2}{}\fi
\noindent\textit{Source reference.} Original lemma 0911: focused-seeds/AGP/sources/stacks-snapshot-20261011/tags.txt, source lines 11429--11429. Navigation pointer to a statement the counted proof refers to but does not contain; nothing load-bearing is delegated to it.\par\medskip
\ifdefined\hypertarget\hypertarget{sourcepointer-3}{}\fi
\noindent\textit{Source reference.} Original lemma 070E: focused-seeds/AGP/sources/stacks-snapshot-20261011/tags.txt, source lines 8953--8953. Navigation pointer to a statement the counted proof refers to but does not contain; nothing load-bearing is delegated to it.\par\medskip
\ifdefined\hypertarget\hypertarget{sourcepointer-4}{}\fi
\noindent\textit{Source reference.} Original lemma 00HD: focused-seeds/AGP/sources/stacks-snapshot-20261011/tags.txt, source lines 639--639. Navigation pointer to a statement the counted proof refers to but does not contain; nothing load-bearing is delegated to it.\par\medskip
\begin{lemma}
\label{lemma-limit-flat}
Let $A$ be a Noetherian ring. Let $I$ be an ideal of $A$. Let
$(M_n)$ be an inverse system of $A$-modules such that
\begin{enumerate}
\item $M_n$ is a flat $A/I^n$-module,
\item $M_{n + 1} \to M_n$ is surjective.
\end{enumerate}
Then $M = \lim M_n$ is a flat $A$-module and
$Q \otimes_A M = \lim Q \otimes_A M_n$ for every finite $A$-module $Q$.
\end{lemma}

\begin{proof}
We first show that $Q \otimes_A M = \lim Q \otimes_A M_n$ for every finite
$A$-module $Q$. Choose a resolution $F_2 \to F_1 \to F_0 \to Q \to 0$
by finite free $A$-modules $F_i$. Then
$$
F_2 \otimes_A M_n \to F_1 \otimes_A M_n \to F_0 \otimes_A M_n
$$
is a chain complex whose homology in degree $0$ is $Q \otimes_A M_n$
and whose homology in degree $1$ is
$$
\text{Tor}_1^A(Q, M_n) = \text{Tor}_1^A(Q, A/I^n) \otimes_{A/I^n} M_n
$$
as $M_n$ is flat over $A/I^n$. By Lemma \ref{lemma-tor-strictly-pro-zero}
we see that this system is essentially constant (with value $0$).
It follows from Homology, Lemma \ref{homology-lemma-apply-Mittag-Leffler-again}
that $\lim Q \otimes_A M_n =
\Coker(\lim F_1 \otimes_A M_n \to \lim F_0 \otimes_A M_n)$.
Since $F_i$ is finite free this equals
$\Coker(F_1 \otimes_A M \to F_0 \otimes_A M) = Q \otimes_A M$.

\medskip\noindent
Next, let $Q \to Q'$ be an injective map of finite $A$-modules.
We have to show that $Q \otimes_A M \to Q' \otimes_A M$ is injective
(Algebra, Lemma \ref{algebra-lemma-flat}). By the above we see
$$
\Ker(Q \otimes_A M \to Q' \otimes_A M) =
\Ker(\lim Q \otimes_A M_n \to \lim Q' \otimes_A M_n).
$$
For each $n$ we have an exact sequence
$$
\text{Tor}_1^A(Q', M_n) \to \text{Tor}_1^A(Q'', M_n) \to
Q \otimes_A M_n \to Q' \otimes_A M_n
$$
where $Q'' = \Coker(Q \to Q')$. Above we have seen that the
inverse systems of Tor's are essentially constant with value $0$.
It follows from
Homology, Lemma \ref{homology-lemma-apply-Mittag-Leffler-again}
that the inverse limit of the right most maps is injective.
\end{proof}

\end{document}
